\documentclass[11pt,a4paper]{article}
\usepackage[utf8]{inputenc}
\usepackage[T1]{fontenc}
\usepackage[spanish,es-minimal,es-noshorthands]{babel}
\AtBeginDocument{\renewcommand{\sin}{\operatorname{sen}}\renewcommand{\arcsin}{\operatorname{arcsen}}}
\usepackage{lmodern}
\usepackage[sfdefault,lining]{FiraSans}
\usepackage[a4paper,top=2.3cm,bottom=2.2cm,left=1.8cm,right=1.8cm,headheight=1.3cm,headsep=0.4cm,footskip=1cm]{geometry}
\usepackage{amsmath,amssymb}
\usepackage{icomma}
\usepackage{graphicx}
\usepackage[table]{xcolor}
\usepackage{tikz}
\usepackage{fontawesome5}
\usepackage{ccicons}
\usepackage{fancyhdr}
\usepackage{lastpage}
\usepackage{tasks}
\usepackage{needspace}
\usepackage[most]{tcolorbox}
\usepackage{array}
\usepackage[hidelinks]{hyperref}
\AtBeginDocument{\hypersetup{pdftitle={Soluciones de la ficha de la unidad 6: Aplicaciones de la derivada. 2º Bach.},pdfsubject={Matemáticas CCSS 2º Bach.},pdfkeywords={aplicaciones de la derivada, recta tangente, recta normal, monotonía, crecimiento, decrecimiento, máximos, mínimos, extremos absolutos, optimización, PAU}}}

\definecolor{navy}{HTML}{1B2748}
\definecolor{accent}{HTML}{E8603A}
\definecolor{teal}{HTML}{1F9E90}
\definecolor{soft}{HTML}{EEF1F8}
\definecolor{muted}{HTML}{5D6685}

\setlength{\parindent}{0pt}
\setlength{\parskip}{3pt}
\settasks{label=\alph*), label-format=\bfseries\color{navy}, label-width=1.3em, item-indent=1.8em, column-sep=0.6em, after-item-skip=5pt, before-skip=3pt, after-skip=4pt}

% ---------- capçalera i peu ----------
\newcommand{\llicencia}{\href{https://creativecommons.org/licenses/by-nc-sa/4.0/deed.es}{\ccbyncsa}\ \ Licencia Creative Commons Atribución-NoComercial-CompartirIgual 4.0 Internacional (CC BY-NC-SA 4.0)}
\pagestyle{fancy}
\fancyhf{}
\fancyhead[L]{\raisebox{-0.35\height}{\includegraphics[height=0.95cm]{logo-gva.png}}}
\fancyhead[C]{\small\color{muted}Soluciones de la ficha · Unidad 6: Aplicaciones de la derivada · 2\textsuperscript{o} Bach.}
\fancyhead[R]{\raisebox{-0.35\height}{\includegraphics[height=0.95cm]{logo-vermellar.png}}}
\fancyfoot[L]{\scriptsize\color{muted}\llicencia}
\fancyfoot[R]{\small\color{muted}\thepage}
\renewcommand{\headrulewidth}{0.4pt}
\renewcommand{\headrule}{\color{muted!40}\hrule height \headrulewidth}
\fancypagestyle{primera}{\fancyhf{}\renewcommand{\headrulewidth}{0pt}
  \fancyfoot[L]{\scriptsize\color{muted}\llicencia}
  \fancyfoot[R]{\small\color{muted}\thepage}}

% ---------- nivells i exercicis ----------
\newcommand{\estrelles}[1]{{\color{accent}\ifcase#1\or\faStar\faStar[regular]\faStar[regular]\or\faStar\faStar\faStar[regular]\or\faStar\faStar\faStar\fi}}
\newcounter{exer}
\newcommand{\ex}[1]{\par\needspace{4\baselineskip}\medskip\refstepcounter{exer}\noindent{\color{navy}\bfseries\theexer.}\ \ }

\newcommand{\seccio}[2]{\par\needspace{10\baselineskip}\bigskip
  \begin{tcolorbox}[colback=navy,colframe=navy,coltext=white,arc=2mm,boxrule=0pt,left=3mm,right=3mm,top=1.2mm,bottom=1.2mm]
  {\large\bfseries #1\quad #2}
  \end{tcolorbox}\vspace{-2pt}}


% quadre de recordatori (només a la fitxa d'exercicis)
\newcommand{\recorda}[1]{\par\smallskip\begin{tcolorbox}[colback=accent!7,colframe=accent!55,arc=2mm,boxrule=0.5pt,left=2.5mm,right=2.5mm,top=1mm,bottom=1mm]\small{\color{accent}\faLightbulb[regular]}\ \textbf{Recuerda:}\ #1\end{tcolorbox}}
% etiqueta dels exercicis de la PAU d'altres comunitats
\newcommand{\pau}[1]{{\color{teal}\small\bfseries\faUniversity\ PAU · #1.}\ }
\newcommand{\res}[1]{\boxed{#1}}\newcommand{\rest}[1]{\fbox{#1}}

% =====================================================================
% Capçalera FentMates de les fitxes: un color per a cada curs.
% Els colors són els de la pàgina inicial (LEVELS a assets/js/materials.js);
% el logo ve de _eines/logo/pdf/ (python3 _eines/logo/genera.py).
%
% Ús, al final del preàmbul de la unitat:
%   % =====================================================================
% Capçalera FentMates de les fitxes: un color per a cada curs.
% Els colors són els de la pàgina inicial (LEVELS a assets/js/materials.js);
% el logo ve de _eines/logo/pdf/ (python3 _eines/logo/genera.py).
%
% Ús, al final del preàmbul de la unitat:
%   % =====================================================================
% Capçalera FentMates de les fitxes: un color per a cada curs.
% Els colors són els de la pàgina inicial (LEVELS a assets/js/materials.js);
% el logo ve de _eines/logo/pdf/ (python3 _eines/logo/genera.py).
%
% Ús, al final del preàmbul de la unitat:
%   \input{../comu/fentmates.tex}
%   \fmcurs{1eso}
% i, al començament del document:
%   \capcalera{Matemàtiques}{1r ESO}{Fitxa d'exercicis · Unitat 2}{Divisibilitat}
% Opcional abans de \input: \def\fmneed{14} (espai mínim abans de cada secció).
% =====================================================================
\makeatletter
\def\fm@logodir{../../logo/pdf}
\providecommand{\fmneed}{10}
\def\fm@cursnom{}

% clau del curs, color i color del text damunt del color
\newcommand{\fm@curs}[3]{\@namedef{fm@color@#1}{#2}\@namedef{fm@text@#1}{#3}}
\fm@curs{1eso}{E6B84A}{navy}
\fm@curs{2eso}{3B82F6}{white}
\fm@curs{3eso}{FF7A50}{white}
\fm@curs{4eso-a}{A78BFA}{white}
\fm@curs{4eso-b}{6D28D9}{white}
\fm@curs{1batx-mi}{9B1C1C}{white}
\fm@curs{1batx-ccss}{EF6B6B}{white}
\fm@curs{2batx-mii}{1E6B47}{white}
\fm@curs{2batx-ccss}{2BB3A3}{white}

\newcommand{\fmcurs}[1]{%
  \@ifundefined{fm@color@#1}{\PackageError{fentmates}{Curs desconegut: #1}{}}{}%
  \edef\fm@tmp{\noexpand\definecolor{curs}{HTML}{\@nameuse{fm@color@#1}}%
               \noexpand\colorlet{curstext}{\@nameuse{fm@text@#1}}}\fm@tmp
  \colorlet{cursfosc}{curs!72!black}%
  \def\fm@logo{\fm@logodir/vertical-#1.pdf}%
  \def\fm@marca{\fm@logodir/marca-#1.pdf}}

% ---------- capçalera de la primera pàgina ----------
% #1 matèria (petit), #2 curs, #3 tipus de document i unitat, #4 títol
\newcommand{\capcalera}[4]{%
  \gdef\fm@cursnom{#2}%
  \par\noindent
  \begin{tcolorbox}[colback=curs,colframe=curs,boxrule=0pt,arc=2.2mm,width=3.5cm,height=2.15cm,
      valign=center,left=3mm,right=2mm,top=1mm,bottom=1mm,nobeforeafter,box align=center]
    \color{curstext}{\footnotesize #1}\par\vspace{1pt}{\fontsize{19}{21}\selectfont\bfseries #2\par}%
  \end{tcolorbox}\hspace{5mm}%
  \begin{minipage}[c]{\dimexpr\linewidth-3.5cm-5mm-2.9cm\relax}
    {\color{cursfosc}\small\bfseries\MakeUppercase{#3}}\par\vspace{2pt}
    {\color{navy}\fontsize{20}{23}\selectfont\bfseries #4\par}
  \end{minipage}\hfill
  \begin{minipage}[c]{2.6cm}\raggedleft\includegraphics[width=2.6cm]{\fm@logo}\end{minipage}\par
  \vspace{2pt}}

% ---------- barres de secció del color del curs ----------
\renewcommand{\seccio}[2]{\par\needspace{\fmneed\baselineskip}\bigskip
  \begin{tcolorbox}[colback=curs,colframe=curs,coltext=curstext,arc=2mm,boxrule=0pt,left=3mm,right=3mm,top=1.2mm,bottom=1.2mm]
  {\large\bfseries #1\quad #2}
  \end{tcolorbox}\vspace{-2pt}}

% ---------- estrelles de dificultat del color del curs (a Batxillerat no se'n posen) ----------
\@ifundefined{estrelles}{}{%
  \renewcommand{\estrelles}[1]{{\color{curs}\ifcase#1\or\faStar\faStar[regular]\faStar[regular]\or\faStar\faStar\faStar[regular]\or\faStar\faStar\faStar\fi}}}

% ---------- capçalera i peu de les pàgines ----------
% primera pàgina: logos institucionals al peu
\fancypagestyle{primera}{\fancyhf{}\renewcommand{\headrulewidth}{0pt}
  \fancyfoot[L]{\raisebox{-0.32\height}{\includegraphics[height=0.78cm]{logo-gva.png}}\hspace{5mm}%
    \raisebox{-0.32\height}{\includegraphics[height=0.78cm]{logo-vermellar.png}}}
  \fancyfoot[R]{\scriptsize\color{muted}\href{https://creativecommons.org/licenses/by-nc-sa/4.0/deed.ca}{\ccbyncsa}\ \ CC BY-NC-SA 4.0 · fentmates.org\qquad\small\thepage}}
% la resta: el símbol a l'esquerra i el curs a la dreta, amb el títol al centre com abans
\fancyhead[L]{\raisebox{-0.3\height}{\includegraphics[height=0.62cm]{\fm@marca}}}
\fancyhead[R]{\small\bfseries\color{cursfosc}\fm@cursnom}
\renewcommand{\headrule}{\color{curs!45}\hrule height \headrulewidth}
\makeatother

%   \fmcurs{1eso}
% i, al començament del document:
%   \capcalera{Matemàtiques}{1r ESO}{Fitxa d'exercicis · Unitat 2}{Divisibilitat}
% Opcional abans de \input: \def\fmneed{14} (espai mínim abans de cada secció).
% =====================================================================
\makeatletter
\def\fm@logodir{../../logo/pdf}
\providecommand{\fmneed}{10}
\def\fm@cursnom{}

% clau del curs, color i color del text damunt del color
\newcommand{\fm@curs}[3]{\@namedef{fm@color@#1}{#2}\@namedef{fm@text@#1}{#3}}
\fm@curs{1eso}{E6B84A}{navy}
\fm@curs{2eso}{3B82F6}{white}
\fm@curs{3eso}{FF7A50}{white}
\fm@curs{4eso-a}{A78BFA}{white}
\fm@curs{4eso-b}{6D28D9}{white}
\fm@curs{1batx-mi}{9B1C1C}{white}
\fm@curs{1batx-ccss}{EF6B6B}{white}
\fm@curs{2batx-mii}{1E6B47}{white}
\fm@curs{2batx-ccss}{2BB3A3}{white}

\newcommand{\fmcurs}[1]{%
  \@ifundefined{fm@color@#1}{\PackageError{fentmates}{Curs desconegut: #1}{}}{}%
  \edef\fm@tmp{\noexpand\definecolor{curs}{HTML}{\@nameuse{fm@color@#1}}%
               \noexpand\colorlet{curstext}{\@nameuse{fm@text@#1}}}\fm@tmp
  \colorlet{cursfosc}{curs!72!black}%
  \def\fm@logo{\fm@logodir/vertical-#1.pdf}%
  \def\fm@marca{\fm@logodir/marca-#1.pdf}}

% ---------- capçalera de la primera pàgina ----------
% #1 matèria (petit), #2 curs, #3 tipus de document i unitat, #4 títol
\newcommand{\capcalera}[4]{%
  \gdef\fm@cursnom{#2}%
  \par\noindent
  \begin{tcolorbox}[colback=curs,colframe=curs,boxrule=0pt,arc=2.2mm,width=3.5cm,height=2.15cm,
      valign=center,left=3mm,right=2mm,top=1mm,bottom=1mm,nobeforeafter,box align=center]
    \color{curstext}{\footnotesize #1}\par\vspace{1pt}{\fontsize{19}{21}\selectfont\bfseries #2\par}%
  \end{tcolorbox}\hspace{5mm}%
  \begin{minipage}[c]{\dimexpr\linewidth-3.5cm-5mm-2.9cm\relax}
    {\color{cursfosc}\small\bfseries\MakeUppercase{#3}}\par\vspace{2pt}
    {\color{navy}\fontsize{20}{23}\selectfont\bfseries #4\par}
  \end{minipage}\hfill
  \begin{minipage}[c]{2.6cm}\raggedleft\includegraphics[width=2.6cm]{\fm@logo}\end{minipage}\par
  \vspace{2pt}}

% ---------- barres de secció del color del curs ----------
\renewcommand{\seccio}[2]{\par\needspace{\fmneed\baselineskip}\bigskip
  \begin{tcolorbox}[colback=curs,colframe=curs,coltext=curstext,arc=2mm,boxrule=0pt,left=3mm,right=3mm,top=1.2mm,bottom=1.2mm]
  {\large\bfseries #1\quad #2}
  \end{tcolorbox}\vspace{-2pt}}

% ---------- estrelles de dificultat del color del curs (a Batxillerat no se'n posen) ----------
\@ifundefined{estrelles}{}{%
  \renewcommand{\estrelles}[1]{{\color{curs}\ifcase#1\or\faStar\faStar[regular]\faStar[regular]\or\faStar\faStar\faStar[regular]\or\faStar\faStar\faStar\fi}}}

% ---------- capçalera i peu de les pàgines ----------
% primera pàgina: logos institucionals al peu
\fancypagestyle{primera}{\fancyhf{}\renewcommand{\headrulewidth}{0pt}
  \fancyfoot[L]{\raisebox{-0.32\height}{\includegraphics[height=0.78cm]{logo-gva.png}}\hspace{5mm}%
    \raisebox{-0.32\height}{\includegraphics[height=0.78cm]{logo-vermellar.png}}}
  \fancyfoot[R]{\scriptsize\color{muted}\href{https://creativecommons.org/licenses/by-nc-sa/4.0/deed.ca}{\ccbyncsa}\ \ CC BY-NC-SA 4.0 · fentmates.org\qquad\small\thepage}}
% la resta: el símbol a l'esquerra i el curs a la dreta, amb el títol al centre com abans
\fancyhead[L]{\raisebox{-0.3\height}{\includegraphics[height=0.62cm]{\fm@marca}}}
\fancyhead[R]{\small\bfseries\color{cursfosc}\fm@cursnom}
\renewcommand{\headrule}{\color{curs!45}\hrule height \headrulewidth}
\makeatother

%   \fmcurs{1eso}
% i, al començament del document:
%   \capcalera{Matemàtiques}{1r ESO}{Fitxa d'exercicis · Unitat 2}{Divisibilitat}
% Opcional abans de \input: \def\fmneed{14} (espai mínim abans de cada secció).
% =====================================================================
\makeatletter
\def\fm@logodir{../../logo/pdf}
\providecommand{\fmneed}{10}
\def\fm@cursnom{}

% clau del curs, color i color del text damunt del color
\newcommand{\fm@curs}[3]{\@namedef{fm@color@#1}{#2}\@namedef{fm@text@#1}{#3}}
\fm@curs{1eso}{E6B84A}{navy}
\fm@curs{2eso}{3B82F6}{white}
\fm@curs{3eso}{FF7A50}{white}
\fm@curs{4eso-a}{A78BFA}{white}
\fm@curs{4eso-b}{6D28D9}{white}
\fm@curs{1batx-mi}{9B1C1C}{white}
\fm@curs{1batx-ccss}{EF6B6B}{white}
\fm@curs{2batx-mii}{1E6B47}{white}
\fm@curs{2batx-ccss}{2BB3A3}{white}

\newcommand{\fmcurs}[1]{%
  \@ifundefined{fm@color@#1}{\PackageError{fentmates}{Curs desconegut: #1}{}}{}%
  \edef\fm@tmp{\noexpand\definecolor{curs}{HTML}{\@nameuse{fm@color@#1}}%
               \noexpand\colorlet{curstext}{\@nameuse{fm@text@#1}}}\fm@tmp
  \colorlet{cursfosc}{curs!72!black}%
  \def\fm@logo{\fm@logodir/vertical-#1.pdf}%
  \def\fm@marca{\fm@logodir/marca-#1.pdf}}

% ---------- capçalera de la primera pàgina ----------
% #1 matèria (petit), #2 curs, #3 tipus de document i unitat, #4 títol
\newcommand{\capcalera}[4]{%
  \gdef\fm@cursnom{#2}%
  \par\noindent
  \begin{tcolorbox}[colback=curs,colframe=curs,boxrule=0pt,arc=2.2mm,width=3.5cm,height=2.15cm,
      valign=center,left=3mm,right=2mm,top=1mm,bottom=1mm,nobeforeafter,box align=center]
    \color{curstext}{\footnotesize #1}\par\vspace{1pt}{\fontsize{19}{21}\selectfont\bfseries #2\par}%
  \end{tcolorbox}\hspace{5mm}%
  \begin{minipage}[c]{\dimexpr\linewidth-3.5cm-5mm-2.9cm\relax}
    {\color{cursfosc}\small\bfseries\MakeUppercase{#3}}\par\vspace{2pt}
    {\color{navy}\fontsize{20}{23}\selectfont\bfseries #4\par}
  \end{minipage}\hfill
  \begin{minipage}[c]{2.6cm}\raggedleft\includegraphics[width=2.6cm]{\fm@logo}\end{minipage}\par
  \vspace{2pt}}

% ---------- barres de secció del color del curs ----------
\renewcommand{\seccio}[2]{\par\needspace{\fmneed\baselineskip}\bigskip
  \begin{tcolorbox}[colback=curs,colframe=curs,coltext=curstext,arc=2mm,boxrule=0pt,left=3mm,right=3mm,top=1.2mm,bottom=1.2mm]
  {\large\bfseries #1\quad #2}
  \end{tcolorbox}\vspace{-2pt}}

% ---------- estrelles de dificultat del color del curs (a Batxillerat no se'n posen) ----------
\@ifundefined{estrelles}{}{%
  \renewcommand{\estrelles}[1]{{\color{curs}\ifcase#1\or\faStar\faStar[regular]\faStar[regular]\or\faStar\faStar\faStar[regular]\or\faStar\faStar\faStar\fi}}}

% ---------- capçalera i peu de les pàgines ----------
% primera pàgina: logos institucionals al peu
\fancypagestyle{primera}{\fancyhf{}\renewcommand{\headrulewidth}{0pt}
  \fancyfoot[L]{\raisebox{-0.32\height}{\includegraphics[height=0.78cm]{logo-gva.png}}\hspace{5mm}%
    \raisebox{-0.32\height}{\includegraphics[height=0.78cm]{logo-vermellar.png}}}
  \fancyfoot[R]{\scriptsize\color{muted}\href{https://creativecommons.org/licenses/by-nc-sa/4.0/deed.ca}{\ccbyncsa}\ \ CC BY-NC-SA 4.0 · fentmates.org\qquad\small\thepage}}
% la resta: el símbol a l'esquerra i el curs a la dreta, amb el títol al centre com abans
\fancyhead[L]{\raisebox{-0.3\height}{\includegraphics[height=0.62cm]{\fm@marca}}}
\fancyhead[R]{\small\bfseries\color{cursfosc}\fm@cursnom}
\renewcommand{\headrule}{\color{curs!45}\hrule height \headrulewidth}
\makeatother

\fmcurs{2batx-ccss}

\begin{document}
\thispagestyle{primera}

\capcalera{Matemáticas CCSS}{2\textsuperscript{o} Bach.}{Soluciones de la ficha · Unidad 6}{Aplicaciones de la derivada}
\vspace{8pt}

\begin{tcolorbox}[colback=soft,colframe=soft,arc=2mm,left=3mm,right=3mm,top=1.5mm,bottom=1.5mm]
\small Intenta resolver cada ejercicio antes de mirar su solución. Si no te sale el mismo resultado, repasa los pasos o compruébalo con la herramienta de estudio de la monotonía y los extremos de los apuntes interactivos de la unidad.
\end{tcolorbox}

% =====================================================
\seccio{1}{Recta tangente y recta normal}

\ex{1} \begin{tasks}(1)
  \task $f(2)=1$, $f'(x)=3x^2-4x$, $f'(2)=4$: \ tangente $\res{y=4x-7}$; \ normal $y-1=-\frac14(x-2)\Rightarrow\res{y=-\frac14x+\frac32}$
  \task $f(0)=-1$, $f'(x)=\frac{3}{(x+1)^2}$, $f'(0)=3$: \ tangente $\res{y=3x-1}$; \ normal $\res{y=-\frac13x-1}$
  \task $f(1)=0$, $f'(x)=\ln x+1$, $f'(1)=1$: \ tangente $\res{y=x-1}$; \ normal $\res{y=-x+1}$
  \task $f(1)=1$, $f'(x)=2x\,e^{x^2-1}$, $f'(1)=2$: \ tangente $\res{y=2x-1}$; \ normal $\res{y=-\frac12x+\frac32}$
\end{tasks}

\ex{2} $f'(x)=3x^2-12x=-9\Rightarrow x^2-4x+3=0\Rightarrow x=1,\ x=3$. \ $f(1)=3$: $\res{y=-9x+12}$; \ $f(3)=-19$: $\res{y=-9x+8}$.

\ex{2} $-\frac{1}{f'(a)}=-3\Rightarrow f'(a)=\frac13$. \ $f'(x)=\frac{1}{\sqrt{2x+1}}=\frac13\Rightarrow2x+1=9\Rightarrow x=4$, punto $\res{(4,3)}$. \ Normal: $y-3=-3(x-4)\Rightarrow\res{y=-3x+15}$.

\ex{1} \pau{Madrid, EvAU julio 2019}$f'(x)=6x^2-8=0\Rightarrow x=\pm\frac{2}{\sqrt3}=\pm\frac{2\sqrt3}{3}$. \ Puntos $\res{\left(\frac{2\sqrt3}{3},-\frac{32\sqrt3}{9}\right)}$ y $\res{\left(-\frac{2\sqrt3}{3},\frac{32\sqrt3}{9}\right)}$.

\ex{2} \pau{Cataluña, PAU junio 2019}Tangente horizontal: $f'(x)=a-\frac{8}{x^2}\Rightarrow f'(-2)=a-2=0\Rightarrow\res{a=2}$. \ $f(-2)=-4+b-4=-6\Rightarrow\res{b=2}$.

\ex{2} \pau{Andalucía, PEvAU junio 2018}\begin{tasks}(1)
  \task Continuidad: $1+a=b+2$. \ Derivabilidad: $3+2a=b-2$. \ $\res{a=-6,\ b=-7}$.
  \task En $x=2$ usamos $f(x)=3x+\frac2x$: $f(2)=7$, $f'(x)=3-\frac{2}{x^2}$, $f'(2)=\frac52$: \ $y-7=\frac52(x-2)\Rightarrow\res{5x-2y+4=0}$.
\end{tasks}

\ex{1} \pau{Murcia, EBAU septiembre 2019}$f(1)=2e^2$, $f'(x)=2e^{x+1}$, $f'(1)=2e^2$: \ $y-2e^2=2e^2(x-1)\Rightarrow\res{y=2e^2x}$.


% =====================================================
\seccio{2}{Monotonía de una función}

\ex{1} \begin{tasks}(1)
  \task $f'(x)=3(x+1)(x-3)$: creciente en $\res{(-\infty,-1)\cup(3,+\infty)}$, decreciente en $\res{(-1,3)}$.
  \task $f'(x)=4x(x+2)(x-2)$: decreciente en $\res{(-\infty,-2)\cup(0,2)}$, creciente en $\res{(-2,0)\cup(2,+\infty)}$.
  \task $f'(x)=-3(x+1)(x-1)$: decreciente en $\res{(-\infty,-1)\cup(1,+\infty)}$, creciente en $\res{(-1,1)}$.
  \task $f'(x)=3x^2+3>0$ siempre: \rest{creciente en todo $\mathbb R$}.
\end{tasks}

\ex{2} \begin{tasks}(1)
  \task $f'(x)=\frac{-2}{(x-2)^2}<0$: decreciente en $\res{(-\infty,2)}$ y en $\res{(2,+\infty)}$.
  \task $f'(x)=\frac{x^2-1}{x^2}$: creciente en $\res{(-\infty,-1)\cup(1,+\infty)}$, decreciente en $\res{(-1,0)\cup(0,1)}$.
  \task $f'(x)=\frac{-2x}{(x^2+1)^2}$: creciente en $\res{(-\infty,0)}$, decreciente en $\res{(0,+\infty)}$.
  \task $f'(x)=\frac{-2x}{(x^2-1)^2}$: creciente en $\res{(-\infty,-1)\cup(-1,0)}$, decreciente en $\res{(0,1)\cup(1,+\infty)}$.
\end{tasks}

\ex{2} \begin{tasks}(1)
  \task $f'(x)=x(x+2)e^x$: creciente en $\res{(-\infty,-2)\cup(0,+\infty)}$, decreciente en $\res{(-2,0)}$.
  \task $f'(x)=\frac{e^x(x-1)}{x^2}$: decreciente en $\res{(-\infty,0)\cup(0,1)}$, creciente en $\res{(1,+\infty)}$.
  \task $f'(x)=\frac{2x}{x^2+4}$: decreciente en $\res{(-\infty,0)}$, creciente en $\res{(0,+\infty)}$.
  \task Dominio $(0,+\infty)$. $f'(x)=2x-\frac2x=\frac{2(x^2-1)}{x}$: decreciente en $\res{(0,1)}$, creciente en $\res{(1,+\infty)}$.
\end{tasks}

\ex{1} \begin{tasks}(1)
  \task Creciente en $\res{(-\infty,-1)\cup(1,+\infty)}$ y decreciente en $\res{(-1,1)}$.
  \task Máximo relativo $\res{(-1,2)}$ y mínimo relativo $\res{(1,-2)}$.
  \task $\res{f'(0)<0}$ (decrece), $\res{f'(2)>0}$ (crece) y $\res{f'(-1)=0}$ (tangente horizontal).
\end{tasks}

\ex{2} \pau{Madrid, EvAU junio 2018}\begin{tasks}(1)
  \task $\res{\mathbb R-\{-1\}}$
  \task $f'(x)=\frac{3x^2(x+1)^2-2x^3(x+1)}{(x+1)^4}=\frac{x^2(x+3)}{(x+1)^3}$. Creciente en $\res{(-\infty,-3)\cup(-1,+\infty)}$ y decreciente en $\res{(-3,-1)}$. (En $x=0$, $f'=0$ pero no cambia de signo.)
\end{tasks}

\ex{2} \pau{Islas Baleares, PBAU ordinaria 2021}\begin{tasks}(1)
  \task $f'(1)=3a+2b+1=0$ y $f(3)=27a+9b+3=0$: \ $\res{a=\frac19,\ b=-\frac23}$.
  \task $f'(x)=9x^2+6x+1=(3x+1)^2\geq0$: \rest{creciente en todo $\mathbb R$} (en $x=-\frac13$ la tangente es horizontal, pero no hay extremo).
\end{tasks}

\ex{2} \pau{Canarias, EBAU junio 2017}\begin{tasks}(1)
  \task $f'(x)=-0{,}5<0$ si $0<x<6$ y $f'(x)=\frac{5(x+4)-5(x+14)}{(x+4)^2}=\frac{-50}{(x+4)^2}<0$ si $x>6$: \rest{siempre decrece}.
  \task $f(6)=10$ y es decreciente: después revela \rest{menos de 10}. $\lim_{x\to+\infty}\frac{5(x+14)}{x+4}=5$ y se acerca por arriba: \rest{nunca menos de 5}.
\end{tasks}


% =====================================================
\seccio{3}{Máximos y mínimos relativos}

\ex{1} \begin{tasks}(1)
  \task $f'(x)=6(x-1)(x-2)$, $f''(x)=12x-18$: \ máximo $\res{(1,1)}$, mínimo $\res{(2,0)}$.
  \task $f'(x)=4x(x^2-1)$, $f''(x)=12x^2-4$: \ máximo $\res{(0,3)}$, mínimos $\res{(-1,2)}$ y $\res{(1,2)}$.
  \task $f'(x)=1-\frac{1}{x^2}$, $f''(x)=\frac{2}{x^3}$: \ máximo $\res{(-1,-2)}$, mínimo $\res{(1,2)}$.
  \task $f'(x)=e^{-x}(1-x)$, $f''(x)=e^{-x}(x-2)$, $f''(1)=-\frac1e<0$: \ máximo $\res{\left(1,\frac1e\right)}$.
\end{tasks}

\ex{2} $f'(x)=12x^3-12x^2=12x^2(x-1)=0\Rightarrow x=0,\ x=1$. \ $f''(x)=36x^2-24x$: $f''(1)=12>0$, mínimo $\res{(1,-1)}$. \ $f''(0)=0$ no decide: $f'<0$ a la izquierda y a la derecha de 0 (no cambia de signo), así que en $x=0$ \rest{no hay extremo}.

\ex{1} \pau{La Rioja, EBAU julio 2017}\begin{tasks}(1)
  \task $f'(x)=6(x+1)(x-2)$: creciente en $\res{(-\infty,-1)\cup(2,+\infty)}$, decreciente en $\res{(-1,2)}$.
  \task Máximo relativo $\res{(-1,0)}$ y mínimo relativo $\res{(2,-27)}$.
\end{tasks}

\ex{1} \pau{Navarra, EvAU ordinaria 2020}$y'=3x(x-2)$: creciente en $\res{(-\infty,0)\cup(2,+\infty)}$, decreciente en $\res{(0,2)}$. Máximo $\res{(0,0)}$ y mínimo $\res{(2,-4)}$.

\ex{1} \pau{País Vasco, EAU junio 2019}\begin{tasks}(1)
  \task $f'(x)=3(x-1)(x-3)$: creciente en $\res{(-\infty,1)\cup(3,+\infty)}$, decreciente en $\res{(1,3)}$.
  \task Máximo relativo $\res{(1,4)}$ y mínimo relativo $\res{(3,0)}$.
\end{tasks}

\ex{2} \pau{Cataluña, PAU ordinaria 2021}\begin{tasks}(1)
  \task $f'(x)=12x^2+2ax$, $f'(-1)=12-2a=0\Rightarrow\res{a=6}$.
  \task $f'(x)=12x(x+2)$: creciente en $\res{(-\infty,-2)\cup(0,+\infty)}$, decreciente en $\res{(-2,0)}$. Máximo $\res{(-2,14)}$ y mínimo $\res{(0,-2)}$.
\end{tasks}

\ex{2} \pau{Murcia, EBAU julio 2020}\begin{tasks}(1)
  \task $f'(0)=b=-1$ y $f'(1)=3+2a+b=0$: \ $\res{a=-1,\ b=-1}$.
  \task $f'(x)=3x^2-4x-4=(3x+2)(x-2)$: creciente en $\res{\left(-\infty,-\frac23\right)\cup(2,+\infty)}$, decreciente en $\res{\left(-\frac23,2\right)}$. Máximo $\res{\left(-\frac23,\frac{67}{27}\right)}$ y mínimo $\res{(2,-7)}$.
\end{tasks}

\ex{2} \pau{País Vasco, EAU julio 2019}$f(0)=c=0$; \ $f(2)=8+4a+2b=-4$; \ $f'(2)=12+4a+b=0$. \ $\res{a=-3,\ b=0,\ c=0}$. \ $f'(x)=3x(x-2)$: máximo $\res{(0,0)}$ y mínimo $\res{(2,-4)}$.

\ex{2} \pau{La Rioja, EBAU ordinaria 2020}\begin{tasks}(1)
  \task $f'(1)=4-2a=0\Rightarrow a=2$; \ $f(1)=1-2+b=0\Rightarrow b=1$: \ $\res{a=2,\ b=1}$.
  \task Pendiente 0: $f'(x)=4x^3-4x=0\Rightarrow x=0,\pm1$: puntos $\res{(0,1)}$, $\res{(-1,0)}$ y $\res{(1,0)}$.
  \task $f(1)=0$ y $f'(1)=0$: $\res{y=0}$.
\end{tasks}

\ex{3} \pau{Murcia, EBAU septiembre 2019}La pendiente es $m(x)=f'(x)=-3x^2+12x-7$. $m'(x)=-6x+12=0\Rightarrow x=2$ y $m''=-6<0$: máximo. \ $f(2)=7$, $f'(2)=5$: punto $\res{(2,7)}$ y tangente $\res{y=5x-3}$.


% =====================================================
\seccio{4}{Extremos absolutos en un intervalo cerrado}

\ex{1} \begin{tasks}(1)
  \task $f'(x)=3(x^2-1)\Rightarrow x=1$. \ $f(0)=0$, $f(1)=-2$, $f(2)=2$: máximo $\res{2}$ en $x=2$ y mínimo $\res{-2}$ en $x=1$.
  \task $f'(x)=2x-6\Rightarrow x=3$. \ $f(-1)=12$, $f(3)=-4$, $f(4)=-3$: máximo $\res{12}$ en $x=-1$ y mínimo $\res{-4}$ en $x=3$.
  \task $f'(x)=1-\frac{4}{x^2}\Rightarrow x=2$ ($-2$ no es del intervalo). \ $f(1)=5$, $f(2)=4$, $f(5)=5{,}8$: máximo $\res{5{,}8}$ en $x=5$ y mínimo $\res{4}$ en $x=2$.
\end{tasks}

\ex{2} \pau{Andalucía, PEvAU septiembre 2018}\begin{tasks}(1)
  \task $c'(t)=3(t-3)(t-7)$. \ $c(0)=10$, $c(3)=91$, $c(7)=59$, $c(12)=334$: el máximo es en $\res{t=12}$, con $\res{334\,000}$ toneladas.
  \task Decrece en $\res{(3,7)}$.
\end{tasks}

\ex{2} \pau{Extremadura, EBAU ordinaria 2020}$G'(t)=6(t-2)(t-7)$. \ $G(0)=60$, $G(2)=136$, $G(7)=11$, $G(8)=28$. \ Gasto máximo a las $\res{2\text{ h}}$ ($\res{136}$ €) y mínimo a las $\res{7\text{ h}}$ ($\res{11}$ €).

\ex{2} \pau{Extremadura, EBAU julio 2019}$P'(t)=-3(t-2)(t-8)$. \ $P(0)=50$, $P(2)=6$, $P(8)=114$, $P(10)=70$. \ Mínimo a las $\res{2\text{ h}}$ ($\res{6}$ kW) y máximo a las $\res{8\text{ h}}$ ($\res{114}$ kW).

\ex{2} \pau{Aragón, EvAU ordinaria 2020}\begin{tasks}(1)
  \task $C(5)=\res{4{,}5}$ €.
  \task $C'(x)=\frac{2x-16}{10}=0\Rightarrow x=8$. \ $C(2)=7{,}2$, $C(8)=3{,}6$, $C(15)=8{,}5$: mínimo con $\res{8\,000}$ unidades ($\res{3{,}6}$ €) y máximo con $\res{15\,000}$ unidades ($\res{8{,}5}$ €).
\end{tasks}

\ex{2} \pau{Navarra, EvAU junio 2019}\begin{tasks}(1)
  \task $f(0)=\res{20\,000}$ €; \ $f(1)=\res{47\,000}$ €.
  \task $f'(x)=-6x+30=0\Rightarrow x=5$, $f''=-6<0$: gastando $\res{5\,000}$ € el beneficio es máximo, $\res{95\,000}$ €.
  \task Crece en $\res{(0,5)}$ y decrece en $\res{(5,8)}$.
  \task $f(0)=20$, $f(8)=68$: el mínimo es sin gastar en publicidad ($\res{x=0}$), $\res{20\,000}$ €.
\end{tasks}

\ex{3} \pau{Andalucía, PEvAU septiembre 2018}$B'(t)=-0{,}08t+2{,}4$ si $0<t<40$ (cero en $t=30$) y $B'(t)=\frac{320}{t^2}>0$ si $40<t<50$. Crece en $\res{(0,30)\cup(40,50)}$ y decrece en $\res{(30,40)}$. \ $B(0)=0$, $B(30)=36$, $B(40)=32$, $B(50)=33{,}6$: máximo en el año $\res{t=30}$, con $\res{36\,000}$ €.

\ex{2} \pau{Canarias, EBAU julio 2019}\begin{tasks}(1)
  \task $f'(x)=\frac12>0$ en $(0,4)$ y $f'(x)=-2x+12$ en $(4,8)$: crece en $\res{(0,6)}$, decrece en $\res{(6,8)}$ y es constante en $(8,12)$.
  \task $f(0)=4$, $f(6)=10$: máximo en el mes $\res{6}$ ($\res{10}$ millones de €) y mínimo en el mes $\res{0}$ ($\res{4}$ millones de €).
\end{tasks}

\ex{2} \pau{Islas Baleares, PBAU julio 2018}\begin{tasks}(1)
  \task $B'(x)=5>0$ en $(0,3)$ y $B'(x)=-2(x-3)<0$ en $(3,8)$: crece en $\res{(0,3)}$ y decrece en $\res{(3,8)}$.
  \task $B(0)=15$, $B(3)=30$, $B(8)=5$: máximo en $x=3$ ($\res{30\,000}$ €) y mínimo en $x=8$ ($\res{5\,000}$ €).
\end{tasks}


% =====================================================
\seccio{5}{Problemas de optimización: conocemos la función}

\ex{1} \pau{Andalucía, PEvAU junio 2018}$f'(x)=2x-6$: el coste disminuye para $\res{0<x<3}$. \ $f''=2>0$: el coste es mínimo con $\res{x=3}$ artículos, y vale $\res{f(3)=31}$.

\ex{1} \pau{Castilla y León, EBAU julio 2019}\begin{tasks}(1)
  \task $f(60)=\res{5}$.
  \task $f'(t)=-4t+120=0\Rightarrow t=30$, $f''=-4<0$: a los $\res{30}$ minutos, con una atención de $\res{1\,805}$.
\end{tasks}

\ex{1} \pau{Murcia, EBAU julio 2020}$B'(x)=-4x+24=0\Rightarrow x=6$, $B''=-4<0$: vendiendo $\res{6}$ ordenadores, beneficio máximo $\res{36}$.

\ex{1} \pau{Islas Baleares, PBAU ordinaria 2021}\begin{tasks}(1)
  \task $B(110)=\res{4\,800}$ €.
  \task $B'(x)=-2x+300=0\Rightarrow x=150$, $B''<0$: $\res{150}$ unidades y $\res{6\,400}$ €.
\end{tasks}

\ex{2} \pau{Islas Baleares, PBAU ordinaria 2020}\begin{tasks}(1)
  \task $Q(q)=\frac{q^2}{50}+8+\frac{40}{q}$: \ $Q(5)=\res{16{,}5}$ €, \ $Q(20)=\res{18}$ €.
  \task $Q'(q)=\frac{q}{25}-\frac{40}{q^2}=0\Rightarrow q^3=1000\Rightarrow q=10$; \ $Q''(q)=\frac{1}{25}+\frac{80}{q^3}>0$: mínimo con $\res{10}$ mesas, coste medio $\res{14}$ €.
\end{tasks}

\ex{2} \pau{Cataluña, PAU ordinaria 2021}\begin{tasks}(1)
  \task $-0{,}2x^2+5x-20\geq0\Leftrightarrow x^2-25x+100\leq0$: entre $\res{5}$ y $\res{20}$ toneladas.
  \task $f(x)=-0{,}2x+5-\frac{20}{x}$, $f'(x)=-0{,}2+\frac{20}{x^2}=0\Rightarrow x=10$, $f''(10)<0$: $\res{10}$ toneladas, $\res{f(10)=1}$ (1\,000 €).
\end{tasks}

\ex{2} \pau{Aragón, EvAU ordinaria 2021}\begin{tasks}(1)
  \task $B(t)<0\Leftrightarrow 2t-6<0$: durante los $\res{3}$ primeros años.
  \task $B'(t)=\frac{14}{(t+4)^2}>0$: es creciente, así que el máximo es en $\res{t=10}$, con $B(10)=1$, es decir, $\res{100\,000}$ €.
  \task $\frac{2t-6}{t+4}=1{,}5\Rightarrow t=\res{24}$ años.
\end{tasks}

\ex{2} \pau{Aragón, EvAU junio 2018}\begin{tasks}(1)
  \task $B(x)>0\Leftrightarrow x^2-9x+18<0$: $\res{x\in(3,6)}$.
  \task $B'(x)=-\frac{9}{x^2}+\frac{36}{x^3}=0\Rightarrow x=4$. \ $B(1)=-10$, $B(4)=0{,}125$, $B(10)=-0{,}28$: precio de $\res{4}$ €, beneficio de $\res{125\,000}$ €.
\end{tasks}

\ex{2} \pau{Galicia, ABAU julio 2019}\begin{tasks}(1)
  \task $P(0)=4{,}75$: $\res{475}$ €. \ $P(t)=4{,}75\Leftrightarrow t^2-4t=0$: a los $\res{4}$ años.
  \task $P'(t)=\frac{-44(2t-4)}{(t^2-4t+16)^2}$: aumenta en $\res{(0,2)}$ y disminuye a partir de $\res{t=2}$. Máximo a los 2 años: $P(2)=\frac{17}{3}$, unos $\res{566{,}67}$ €.
  \task $\lim_{t\to+\infty}P(t)=2$: tiende a $\res{200}$ €.
\end{tasks}

\ex{2} \pau{Asturias, EBAU julio 2019}\begin{tasks}(1)
  \task $-x^2+6x+3=x+9\Rightarrow x^2-5x+6=0$: a las $\res{2}$ h (11 unidades) y a las $\res{3}$ h (12 unidades).
  \task $f'(x)=-2x+6<0$ si $x>3$: decrece en $\res{(3,6)}$.
  \task $h(x)=f(x)+g(x)=-x^2+7x+12$, $h'(x)=-2x+7=0$: a las $\res{3{,}5}$ h.
\end{tasks}


% =====================================================
\seccio{6}{Problemas de optimización: construimos la función}

\ex{1} \pau{Castilla y León, EBAU septiembre 2017}$y=10-x$, \ $P(x)=x^2(10-x)^2$. \ $P'(x)=2x(10-x)^2-2x^2(10-x)=2x(10-x)(10-2x)=0\Rightarrow x=5$ (los otros valores hacen $P=0$). Los números son $\res{5}$ y $\res{5}$.

\ex{1} \pau{Cataluña, PAU ordinaria 2021}\begin{tasks}(1)
  \task $2x+2y=40\Rightarrow y=20-x$: \ $\res{A(x)=20x-x^2}$.
  \task $A'(x)=20-2x=0\Rightarrow x=10$, $A''<0$: un cuadrado de $\res{10\times10}$ m, con $\res{100}$ m².
\end{tasks}

\ex{2} \pau{La Rioja, EBAU julio 2017}\begin{tasks}(1)
  \task $x$: lados perpendiculares, $y$: lado paralelo. $xy=3200$, $L(x)=2x+\frac{3200}{x}$, $L'(x)=2-\frac{3200}{x^2}=0\Rightarrow x=40$, $L''>0$: $\res{40}$ m perpendicular y $\res{80}$ m paralelo al río.
  \task $2\cdot40\cdot500+80\cdot400=\res{72\,000}$ €.
\end{tasks}

\ex{3} \pau{La Rioja, EBAU junio 2018}\begin{tasks}(1)
  \task $x=\frac{4}{y^2}$, $P(y)=\frac{8}{y^2}+2y$, $P'(y)=-\frac{16}{y^3}+2=0\Rightarrow y=2$: $\res{x=1,\ y=2}$.
  \task $x=1-3y^2$, $A(y)=y-3y^3$, $A'(y)=1-9y^2=0\Rightarrow y=\frac13$: $\res{x=\frac23,\ y=\frac13}$.
\end{tasks}

\ex{2} \pau{Murcia, EBAU junio 2019}$B(x)=40x-(2x^2+4x+98)=-2x^2+36x-98$, $B'(x)=-4x+36=0\Rightarrow x=9$, $B''<0$: $\res{9}$ unidades y $\res{64}$ €.

\ex{2} \pau{Murcia, EBAU junio 2018}$B(x)=15x-C(x)=-2x^2+60x-100$, $B'(x)=-4x+60=0\Rightarrow x=15$, $B''<0$: $\res{15}$ unidades y $\res{350}$ u.m.

\ex{2} \pau{Murcia, EBAU ordinaria 2021}\begin{tasks}(1)
  \task $B(x)=x\left(50-\frac x4\right)-\frac{x^2}{4}-35x-25=-\frac{x^2}{2}+15x-25$, $B'(x)=-x+15=0$: $\res{15}$ unidades.
  \task $p=50-\frac{15}{4}=\res{46{,}25}$ €.
  \task $B(15)=\res{87{,}5}$ €.
\end{tasks}

\ex{2} \pau{Galicia, ABAU extraordinaria 2021}\begin{tasks}(1)
  \task $\res{I(x)=518x-x^3}$, \ $\res{B(x)=-x^3+243x-225}$; \ $B(10)=\res{1\,205}$ €.
  \task $B'(x)=-3x^2+243=0\Rightarrow x=9$, $B''(9)=-54<0$: $\res{9}$ unidades, beneficio $\res{1\,233}$ €, precio $\res{437}$ €.
\end{tasks}

\ex{3} \pau{Cataluña, PAU ordinaria 2020}\begin{tasks}(1)
  \task $B(x)=(18+x-8)(120-4x)=\res{-4x^2+80x+1200}$.
  \task $B'(x)=-8x+80=0\Rightarrow x=10$: hay que \rest{subirlo 10 €}; precio $\res{28}$ € y beneficio $\res{1\,600}$ €.
\end{tasks}

\ex{2} $x$: lado de la base, $h$: altura. $x^2h=4\Rightarrow h=\frac{4}{x^2}$. \ $S(x)=x^2+4xh=x^2+\frac{16}{x}$, $S'(x)=2x-\frac{16}{x^2}=0\Rightarrow x=2$, $S''(2)=6>0$: base de $\res{2\times2}$ m y altura $\res{1}$ m ($12$ m² de material).

\ex{3} $V(x)=x(16-2x)(10-2x)=4x^3-52x^2+160x$, $0<x<5$. \ $V'(x)=12x^2-104x+160=4(3x-20)(x-2)=0\Rightarrow x=2$ ($\frac{20}{3}$ no es válido). $V''(2)=-56<0$: $\res{x=2}$ cm y $\res{V=144}$ cm³ (caja de $12\times6\times2$).


\end{document}
